\scanpage{230}
\mspage{104}
% archive sheet number: 229

Il reste à vérifier (8.24.), qui a été vu dans $\overline{\underline{\unsure{\mathrm{XII}|\mathrm{XV}|\mathrm{XIV}}}}$\ldots\quad / On conclut donc de 7.4. $1^\circ,\unsure{1|2}^\circ,3^\circ$

\uline{Corollaire} 8.3. Posons \unsure{\S|$\xi$}

\[
 \left\{\begin{aligned}
 C_{\deleted{\unsure{A}}}&=\mathrm{Mon}(A^\wedge),&
 \mathrm{TC}_{\deleted{\unsure{A}}}&=C\cap W_{\unsure{A|a}}\\
 F&=(\mathrm{TC})_\ast,&\mathrm{TF}&=C_\ast.
 \end{aligned}\right.
 \tag*{(8.3.1\unsure{.})}
\]

Alors $(C,\mathrm{TF})$ et $(\mathrm{TC},F)$ sont des couples
de Quillen, à factorisation\unsure{s} fonctorielle\unsure{s} (NB
$C$ et $\mathrm{TC}$ \deleted{\unsure{étant}} \inserted{de plus} Q-engendrées \unsure{à p.|i.e.} par des
petites parties de $\mathrm{Fl}\,\widehat A$). Pour que
$(W_A,C,F)$ soit un triple de Quillen \unsure{clos},
il \unsure{f.} et s. qu'on ait

\[
 \underset{\substack{\shortparallel\,\mathrm{def}\\(\mathrm{Mon})_\ast}}{\mathrm{TF}}\subset W_A\ ,
 \tag*{(8.3.2)}
\]

\deleted{\unsure{en}} i.e. \deleted{\unsure{c'est à|on ait}} que toute flèche $f:X\to Y$ \inserted{de type injectif}
de $A^\wedge$, \inserted{(i.e. qui fait de $X$ un objet injectif de $A^\wedge/Y$)} soit dans $W_A$.
Quand cette condition est satisfaite, le triple de Quillen $(W_A,C,F)$ est propre à gauche, par 8.2.1.
Il n'est pas clair qu'elle soit forcément propre -- totalement, i.e. qu'on ait

\[
 F\subset\mathrm{Fib}_{W}\quad (?)\ .
 \tag*{(8.3.3)}
\]

Nous y reviendrons par la suite. Pour l'instant, la chose essentielle, c'est d'examiner la condition (8.3.2) sur $A$.
